% Replacement for the Richardson optimization paragraph in
% `Generalization error (decoder).tex`.

\subsection{HeavyBall optimization term in the decoder generalization bound}

Let
\[
 H_m=G_m^\top G_m+\lambda I,
 \qquad
 z_\lambda^*=H_m^{-1}G_m^\top b_m,
\]
and consider the preconditioned HeavyBall decoder
\begin{equation}
 z_{-1}=z_0=0,
 \qquad
 z_{\ell+1}=z_\ell+B_\theta\alpha
 \bigl(G_m^\top(b_m-G_mz_\ell)-\lambda z_\ell\bigr)
 +\beta(z_\ell-z_{\ell-1}).
 \label{eq:gen-hb}
\end{equation}
Assume $B_\theta\succ0$, $\kappa(B_\theta)\leq\bar\kappa_B$, and
\[
 \sigma(B_\theta^{1/2}H_mB_\theta^{1/2})\subseteq[\mu,M].
\]
With the optimal quadratic HeavyBall coefficients, define
\[
 q=\frac{\sqrt{M/\mu}-1}{\sqrt{M/\mu}+1},
 \qquad
 C_L(q)=1+L(1+q).
\]
The HeavyBall polynomial estimate gives
\begin{equation}
 \|z_L-z_\lambda^*\|_2
 \leq
 \sqrt{\bar\kappa_B}\,C_L(q)q^L
 \|z_0-z_\lambda^*\|_2.
 \label{eq:gen-hb-error}
\end{equation}
Consequently,
\begin{equation}
 \mathbb E\|z_L-z_\lambda^*\|_2^4
 \leq
 \bar\kappa_B^2 C_L(q)^4q^{4L}
 \mathbb E\|z_0-z_\lambda^*\|_2^4.
 \label{eq:gen-hb-fourth}
\end{equation}
Repeating the Galerkin/ridge/optimization decomposition yields
\begin{equation}
\boxed{
\begin{aligned}
 \mathbb E\|\widehat u-u^*\|_{\mathcal H}^4
 \leq{}&C\|D_m\|^4
 \bigl[
 \bar\kappa_B^2 C_L(q)^4q^{4L}
 \mathbb E\|z_0-z_\lambda^*\|_2^4
 +\|z_\lambda^*-z_m^*\|_2^4
 \bigr]\\
 &+C\|u_m^*-u^*\|_{\mathcal H}^4.
\end{aligned}}
\label{eq:gen-hb-fourth-output}
\end{equation}
For $z_0=0$, the first expectation on the right is
$\mathbb E\|z_\lambda^*\|_2^4$.

Thus the Gaussian-tail truncation argument and the $N^{-1/2}$ empirical-process
rate are unchanged.  Only the finite-depth optimization envelope changes from
$q_{\rm Rich}^{4L}$ to
$\bar\kappa_B^2C_L(q)^4q^{4L}$.  The exact HeavyBall polynomial should be used
instead of this envelope when evaluating constants numerically.

\paragraph{Fixed versus selected solver coefficients.}
If $(\alpha,\beta)$ are fixed before the $N$ training tasks are sampled, the
existing one-parameter uniform bound over $\lambda$ applies without a new
parameter-complexity term.  If the same data are used to select
$(\alpha,\beta)$, the empirical process must instead be uniform over a compact
stable set
\[
 \Theta_{\rm HB}
 =\{(\lambda,\alpha,\beta):
 \lambda\in[\lambda_{\min},\lambda_{\max}],
 \qquad 0\leq\beta\leq1-\varepsilon,
 \alpha M\leq2(1+\beta)-\varepsilon\}.
\]
This adds only two scalar coordinates to the covering argument.  Under a
uniform Lipschitz bound for the truncated loss on $\Theta_{\rm HB}$, its metric
entropy is logarithmic in the covering resolution, so the leading
$N^{-1/2}$ rate remains unchanged (up to a logarithmic factor and a
depth-dependent Lipschitz constant).  This statement does not cover an
unrestricted prompt-dependent neural preconditioner; that class requires its
own norm or parameter-space complexity control.

\paragraph{Replica/random-matrix substitution.}
Whenever the Richardson calculation contains the spectral filter
$(1-\eta\lambda)^{2L}$, the HeavyBall calculation replaces it by
\begin{equation}
 p_L(\lambda)^2,
 \quad
 p_{-1}=p_0=1,
 \quad
 p_{\ell+1}=(1+\beta-\alpha\lambda)p_\ell-\beta p_{\ell-1}.
 \label{eq:gen-hb-replica-filter}
\end{equation}
No additional right-hand-side order parameter is needed because
$(\alpha,\beta)$ are fixed across tasks and $p_L$ is independent of the query
$f$.  This is precisely the simplification that fails for PCG.
